\chapter{Paciência}

\lettrine{A}{\space serenidade e a tranquilidade} podem ser extremamente entediantes e trazer à
tona muitas inquietações e dúvidas. A inquietação é um problema comum porque o
reino sensorial é um reino inquieto; os corpos estão inquietos e as mentes estão
inquietas. As condições mudam a todo o momento, portanto se somos apanhados a
reagir às mudanças, ficamos simplesmente inquietos.

É necessário compreender totalmente a inquietação por aquilo que é; a prática
não consiste apenas em usar a vontade para nos prendermos ao tapete da
meditação. Não é um teste para nos tornar numa pessoa forte que tem que vencer a
inquietação -- essa atitude apenas reforça outra visão egoísta. Mas é uma
questão de investigar realmente a inquietação, percebê"-la e conhecê"-la pelo que
é. Para isso, temos que desenvolver paciência; é algo que temos que aprender e
realmente trabalhar.

Quando fui pela primeira vez a Wat Pah Pong, não conseguia entender Laosiano. E
naqueles dias Ajahn Chah estava no seu auge e dando \emph{desanas} (palestras)
de três horas todas as noites. Ele falar podia indefinidamente e todos o amavam
-- ele era um orador muito bom, muito bem"-humorado e todos gostavam das suas
palestras. Mas se não conseguíssemos entender Laosiano\ldots{}!

Eu ficava ali sentado a pensar: ``Quando é que ele vai parar, estou a
desperdiçar o meu tempo.'' Eu ficava mesmo irritado, pensando, ``Estou farto
disto, vou"-me embora.'' Mas eu não conseguia ter coragem suficiente para sair,
então permanecia apenas sentado a pensar -- ``Vou para outro mosteiro. Já chega
disto; não aguento mais.'' E então ele olhava para mim -- ele tinha o sorriso
mais radiante -- e dizia: ``Você está bem?'' E de repente toda a raiva que se
tinha acumulado ao longo daquelas três horas desaparecia completamente.

É curioso, não é? Depois de estar sentado a fumegar por três horas, pode
simplesmente desaparecer. Então, comprometi"-me a que a minha prática seria a
paciência e que durante esse tempo eu desenvolveria a paciência. Eu iria a todas
as palestras e assistiria a todas elas até que o meu físico suportasse.
Determinei"-me a não falhar ou tentar esquivar"-me de alguma, e apenas praticar a
paciência.

E ao fazer isso, comecei a descobrir que a oportunidade de ser paciente foi algo
que me ajudou imenso. A paciência é um alicerce muito firme para a minha visão e
compreensão clara do Dhamma; sem esse alicerce eu teria apenas vagueado e andado
à deriva, como vemos fazer tantas pessoas. Muitos ocidentais chegaram a Wat Pah
Pong e afastaram"-se porque não eram pacientes. Eles não queriam sentar"-se
durante \emph{desanas} de três horas e serem pacientes. Queriam ir a sítios onde
pudessem obter iluminação instantânea e de uma forma rápida, da maneira que
quisessem.

Nós não conseguimos apreciar realmente como a realidade é, através dos desejos e
ambições egoístas que poderão dirigir"-nos, até mesmo no caminho espiritual.
Quando reflecti e realmente contemplei a minha vida em Wat Pah Pong, percebi que
era uma situação muito boa: havia um bom professor, havia o suficiente para
comer, os monges eram bons monges, os leigos eram muito generosos e amáveis e
houve encorajamento para a prática do Dhamma. Melhor do que isso seria difícil;
foi uma oportunidade maravilhosa. E, no entanto, muitos ocidentais não
conseguiam ver isso porque tinham a tendência para pensar -- ``Eu não gosto
disto, eu não quero aquilo; e deveria ser diferente e o que \emph{eu} acho e o
que \emph{eu} sinto, eu não me quero chatear com \emph{isto} e com
\emph{aquilo}.''

Lembro"-me de ir para o mosteiro de Tam Sang Phet, que era um lugar muito
tranquilo e isolado naquela época, morar numa caverna. Um aldeão construiu"-me
uma plataforma porque na parte de baixo da caverna havia uma grande cobra píton.
Uma noite, eu estava sentado nesta plataforma à luz de velas, estava
verdadeiramente arrepiante e a luz projectou sombras em todas as rochas; foi
estranho. Eu estava ali sentado e comecei a ficar muito assustado e, de repente,
fiquei sobressaltado. Olhei para cima e havia uma enorme coruja logo por cima,
olhando para mim. Parecia gigante -- não sei se era assim tão grande, mas
parecia realmente enorme à luz de velas -- e estava a olhar para mim
directamente. Eu pensei, ``Bem, o que há para realmente para ter medo aqui?'' E
tentei imaginar esqueletos e fantasmas ou a Deusa Kali da morte, com presas e
sangue a escorrer da sua boca ou monstros enormes com pele verde e então comecei
a rir porque se tornou muito divertido! Percebi que não estava nem um pouco
assustado.

Naquela época, eu era apenas um monge júnior e uma noite Ajahn Chah levou"-nos a
uma festa de aldeia. Acho que Satimanto Bhikkhu estava lá na época. Éramos todos
praticantes muito sérios e não queríamos qualquer tipo de frivolidade ou
parvoíce; e, claro, ir a uma festa de aldeia seria a última coisa que queríamos
fazer, porque nessas aldeias eles adoram altifalantes. De qualquer forma, Ajahn
Chah levou"-nos, ao Satimanto e a mim a esta festa, e tivemos que ficar sentados
a noite inteira com os sons estridentes dos altifalantes -- e monges a dar
palestras pela noite fora! Não parava de pensar, ``Oh, quero voltar para minha
caverna -- monstros de pele verde e fantasmas são muito melhores do que isto''.
Notei que Satimanto, que estava incrivelmente sério, parecia muito zangado,
crítico e infeliz e nós sentados ali parecíamos simplesmente miseráveis. Eu
pensei: ``Por que é que Ajahn Chah nos traz a estas coisas?''

Então comecei a ver por mim mesmo. Lembro"-me de estar lá sentado a pensar:
``Aqui estou, a ficar todo irritado por causa disto. Isto é assim tão mau? O que
é realmente mau é o que eu estou a fazer disto, o que é realmente miserável é a
minha mente. Altifalantes e barulho, distracção e sonolência, isso nós
conseguimos suportar, mas aquela coisa horrível na minha mente que odeia, que se
ressente e que quer ir embora -- essa é a verdadeira miséria!''

Naquela noite, vi a miséria que poderia criar na minha mente por causa de coisas
que realmente poderia suportar. Lembro"-me disso como uma compreensão claríssima
do que eu pensava ser miserável, e daquilo que realmente é infeliz. No início eu
estava a culpar as pessoas, os altifalantes, a perturbação, o barulho e o
desconforto -- pensei que fosse isso o problema. Então percebi que não era; a
minha mente é que estava miserável.

Se reflectirmos e contemplarmos o Dhamma, aprenderemos com as próprias
situações, das quais menos gostamos -- se tivermos a vontade para o fazer.

\photoPage{image-007-siladhara-bowl.jpg}{Sīladhara a secar a tampa da tigela da mendicância}{Foto por John Baxter}
